\section{Discussion}
\label{chapter:04}

% block-id: 04.discussion.heading
\hypertarget{04.discussion.heading}{}
\subsection{Interpretation of current sharing and loss}

% block-id: 04.discussion.p1
\hypertarget{04.discussion.p1}{}
The three result levels answer the questions posed in the introduction. Round-wire analysis and two-dimensional calibration establish a local skin-effect reference. The 640-strand results show that common terminal voltage permits substantially different complex strand currents, so strand size and interstrand uniformity require separate assessment. Longitudinal PEEC quantifies length and integration effects for its chosen basis. Structural comparisons consequently require control of actual copper area, return, terminals, and axial length, together with transposition-induced changes in coupling. Small discretization changes help identify numerically stable intervals; structural advantages must still be compared with relevant model errors and experimental uncertainty.

% block-id: 04.discussion.p2
\hypertarget{04.discussion.p2}{}
DTM and iDTM \cite{P01,N01} provide a foundation for connecting longitudinal positions with cross-sectional fields. The common-terminal criterion expresses uniform-current requirements as an algebraic condition on the coupled matrix and induced terms, helping test whether a transposition approximation retains important differences. The segmented circuits of iSEC \cite{P02} and current constraints of quasi-two-dimensional methods \cite{P04} also show the need for joint field–circuit solution. Continuum and complex-permeability methods \cite{P05,P09} lower macroscopic computational cost; strand currents and local losses require recovery consistent with their averaging assumptions. These works inform the methods, while our quantitative verification uses its own operating conditions, discretizations, and error records. Comparing experimental accuracy across methods requires a common case.

% block-id: 04.discussion.p3
\hypertarget{04.discussion.p3}{}
The main current uncertainty arises from differences between actual structures and simplified models. A straight two-dimensional bundle omits longitudinal transposition, and longitudinal PEEC omits transverse loops and actual end-current redistribution. The RT0 imposed-field case tests absorption in a fixed cuboid. Supplementary three-dimensional A–φ cylinder calibration retains failures in φ and gauge residuals, while the native three-dimensional Ansys interface has not produced an acceptable field solution. Temperature rise, insulation thickness, winding tolerances, and terminal contacts require thermal modeling and experimental assessment. These factors limit transfer of the present results to actual Litz performance.

% block-id: 04.recommendations.heading
\hypertarget{04.recommendations.heading}{}
\subsection{Priorities for subsequent verification}

% block-id: 04.recommendations.p1
\hypertarget{04.recommendations.p1}{}
Subsequent work should first establish consistent boundaries, gauge treatment, terminal fields, and independent power integrals on manageable three-dimensional cylinders and small-strand structures. Actual terminals and transverse degrees of freedom can then be introduced before comparing traditional multilevel and cross-layer structures meeting the copper-area requirement. Experiments can measure complex impedance at matched frequency, temperature, return, and axial length while recording terminal resistance and geometric tolerances. This sequence separates numerical implementation errors, geometric approximations, and manufacturing deviations and supplies a comparable basis for structural ranking.
